\chapter{Допуск динамического носителя источника центральной щели}
\label{ch:v8-baryon-c0-singlet-triplet-central-gap-dynamical-source-carrier-admission-gate}

Текущий родитель не содержит линейного источника. Проверим, может ли одна
новая динамическая вещественная координата породить его как математическое
ожидание, не вводя внешний член $-j\lambda$.

\section{Минимальный тип носителя}

Пусть $s$ является вещественным семейным и калибровочным синглетом с
чётной градуировкой:
\begin{equation}
 s\in\mathbb R,
 \qquad \rho_G(g)s=s,
 \qquad \rho_{SO(3)_{\rm fam}}(R)s=s,
 \qquad \Gamma_s s=s,
 \qquad \overline s=s.
 \label{eq:v8-baryon-c0-dynamic-source-carrier-type}
\end{equation}
Рассмотрим функционал без внешнего линейного источника
\begin{equation}
 V(\lambda,s)=\frac M2\lambda^2-gs\lambda+\frac u4s^4,
 \qquad M>0,
 \qquad u>0,
 \qquad g\ne0.
 \label{eq:v8-baryon-c0-dynamic-source-parent}
\end{equation}
Исключение $\lambda$ завершением квадрата даёт
\begin{equation}
 V=\frac M2\left(\lambda-\frac gM s\right)^2
 +V_{\rm eff}(s),
 \qquad
 V_{\rm eff}(s)=\frac u4s^4-\frac{g^2}{2M}s^2.
 \label{eq:v8-baryon-c0-dynamic-source-effective-potential}
\end{equation}
Таким образом, смешанный член сам индуцирует отрицательную квадратичную
кривизну для $s$, а квартичный член ограничивает потенциал снизу.

\section{Запуск из нуля}

Гессиан в начале координат равен
\begin{equation}
 H_0=\begin{pmatrix}M&-g\\-g&0\end{pmatrix}.
 \label{eq:v8-baryon-c0-dynamic-source-origin-hessian}
\end{equation}
Его определитель строго отрицателен:
\begin{equation}
 \det H_0=-g^2<0.
 \label{eq:v8-baryon-c0-dynamic-source-origin-instability}
\end{equation}
Следовательно, нулевая конфигурация является седлом без ручного
тахионного массового члена. Уравнения стационарности имеют вид
\begin{equation}
 M\lambda-gs=0,
 \qquad
 us^3-g\lambda=0.
 \label{eq:v8-baryon-c0-dynamic-source-stationarity}
\end{equation}
Кроме нуля существуют две точки, задаваемые точными инвариантами
\begin{equation}
 s_*^2=\frac{g^2}{uM},
 \qquad
 \lambda_* =\frac gM s_*,
 \qquad
 \lambda_*^2=\frac{g^4}{uM^3}.
 \label{eq:v8-baryon-c0-dynamic-source-nonzero-vacua}
\end{equation}
В них индуцированный источник равен $j_*=gs_*$, а энергия понижена:
\begin{equation}
 j_*=gs_*,
 \qquad
 V(\lambda_*,s_*)=-\frac{g^4}{4uM^2}<0.
 \label{eq:v8-baryon-c0-dynamic-source-induced-source-energy}
\end{equation}

Гессиан ненулевого вакуума равен
\begin{equation}
 H_*=\begin{pmatrix}
 M&-g\\
 -g&3g^2/M
 \end{pmatrix}.
 \label{eq:v8-baryon-c0-dynamic-source-vacuum-hessian}
\end{equation}
Критерий Сильвестра даёт
\begin{equation}
 M>0,
 \qquad \det H_*=2g^2>0,
 \qquad H_*>0.
 \label{eq:v8-baryon-c0-dynamic-source-vacuum-stability}
\end{equation}

\section{Операторная совместимость и парность}

Для $h^0=\lambda Q$ и $\Gamma_4^0=2Q$ смешанный инвариант имеет вид
\begin{equation}
 s\Tr(\Gamma_4^0h^0)=\frac32s\lambda.
 \label{eq:v8-baryon-c0-dynamic-source-marked-bilinear}
\end{equation}
Он сохраняет калибровочную группу, семейное $SO(3)$, grading и Real.
Однако полный потенциал сохраняет одновременное отражение
\begin{equation}
 V(-\lambda,-s)=V(\lambda,s).
 \label{eq:v8-baryon-c0-dynamic-source-paired-symmetry}
\end{equation}
Поэтому динамический носитель создаёт ненулевую щель, но не выбирает один
из двух знаков.

\section{Минимальность по удалению членов}

Каждый из трёх членов необходим в данной архитектуре:
\begin{equation}
 \begin{array}{c|c}
 M=0 & \text{неограниченность по }\lambda\text{ при }s\ne0,\\
 g=0 & \lambda_*=s_*=0,\\
 u=0 & V_{\rm eff}(s)=-g^2s^2/(2M)\text{ не ограничен снизу}.
 \end{array}
 \label{eq:v8-baryon-c0-dynamic-source-term-minimality}
\end{equation}
Одномерный вещественный носитель минимален по размерности: без новой
координаты математическое ожидание, заменяющее $j$, отсутствует.

Но нормировка $s$ остаётся координатной. При $s_{\rm new}=cs$, $c>0$,
\begin{equation}
 g_{\rm new}=\frac gc,
 \qquad
 u_{\rm new}=\frac{u}{c^4},
 \qquad
 \frac{g_{\rm new}^4}{u_{\rm new}M^3}
 =\frac{g^4}{uM^3}.
 \label{eq:v8-baryon-c0-dynamic-source-normalization-orbit}
\end{equation}
Значение щели инвариантно, но сами коэффициенты не определены без метрики
носителя.

Итоговые реестры равны
\begin{equation}
 N_{\rm architecture}=8/8,
 \qquad
 N_{\rm origin}=0/4\quad\{\mathbb Rs,M,g,u\}.
 \label{eq:v8-baryon-c0-dynamic-source-ledgers}
\end{equation}

\begin{theorem}[Условный динамический запуск центральной щели]
Функционал $M\lambda^2/2-gs\lambda+us^4/4$ при $M,u>0$ и $g\ne0$
ограничен снизу, имеет седло в нуле и ровно два ненулевых строгих минимума.
В них внешний источник заменяется динамическим значением $j_*=gs_*$.
Архитектура совместима с требуемыми структурными симметриями.
\end{theorem}

\begin{nogocriterion}[Динамический источник не выбирает ориентацию]
Парные вакуумы $(s_*,\lambda_*)$ и $(-s_*,-\lambda_*)$ имеют одинаковую
энергию. Нельзя объявлять знак щели или конкретный носитель выведенным до
идентификации поля $s$, его следовой метрики и члена, снимающего либо
физически интерпретирующего эту парность.
\end{nogocriterion}

\begin{protocol}\begin{enumerate}
\item минимальная размерность носителя: \textbf{$1$ вещественная};
\item внешний линейный источник нужен: \textbf{нет};
\item нуль является седлом: \textbf{да};
\item ненулевых строгих минимумов: \textbf{$2$};
\item индуцированный источник: \textbf{$j_*=gs_*$};
\item знак щели выбран: \textbf{нет};
\item условный архитектурный допуск: \textbf{$8/8$};
\item происхождение носителя и коэффициентов: \textbf{$0/4$};
\item следующий гейт: \textbf{классификация существующих scalar-носителей}.
\end{enumerate}\end{protocol}

Машинный сертификат сохранён в
\path{s2t_v8_baryon_c0_singlet_triplet_central_gap_dynamical_source_carrier_admission_gate_results.json}.